Prolonged sitting is common in occupational and daily life, but total sitting duration alone does not describe whether a seated bout remains static or includes meaningful postural variation. This thesis developed and technically evaluated a low-cost closed-loop sacrum/lower-back wearable system for pelvic-region movement monitoring and detection-outcome-based haptic feedback during prolonged sitting. The prototype combined an XIAO nRF54L15 Sense IMU node, BLE streaming, Android labelled field logging, PC-based signal inspection, compact embedded detection logic, a DRV2605L haptic driver, and a coin vibration motor. Its IMU signals were compared with a research-grade Shimmer IMU in co-located and limited field-worn recordings. A five-participant pilot field study collected labelled seated data for static sitting and instructed anterior-posterior pelvic rotation, which were used for offline classifier evaluation and compact KNN deployment. Co-located Shimmer/XIAO correlations ranged from 0.856 to 0.923 for acceleration axes and from 0.816 to 0.864 for gyroscope axes, while the two usable field-worn paired recordings showed acceleration-motion-magnitude correlations of 0.918 and 0.855 and gyroscope-magnitude correlations of 0.908 and 0.864. On held-out participant P5, the selected full KNN achieved an F1-score of 0.957 for pelvic-rotation detection, and the compact 12-feature, 128-prototype KNN achieved an F1-score of 0.959. Firmware timing measurements showed a mean online decision time of 2.73 ms per 2 s window, with 56.3 KB flash use and 9.3 KB RAM use, and the firmware executed a 10 min inactivity timer with a 1 s local vibration cue. These results support the technical feasibility of the sensing-to-cue pipeline, but they do not establish clinical pelvic kinematics, spontaneous daily-use robustness, long-term wearing comfort, battery lifetime, or behavioural effectiveness of the haptic reminder.
